\documentclass[compress]{beamer}
%,hyperref={pdfpagelabels=false}

\usepackage[ngerman,english]{babel}
\usepackage[T1]{fontenc}
\usepackage[latin1]{inputenc}
\usepackage{amsmath}
\usepackage{amsthm}
\usepackage{amsfonts}
\usepackage{helvet}
\usepackage{url}
\usepackage{listings}
\usepackage{xcolor}
\usepackage{xspace} % Abstand hinter Variablennamen
\usepackage{fix-cm}
%\usepackage[square, sort, numbers, authoryear]{natbib}

\usepackage{beamerthemeLEA2}

%\bibliographystyle{plainnat}

\newcommand{\N}       {\mathbb{N}}          % nat�rliche Zahlen
\newcommand{\Z}       {\mathbb{Z}}          % ganze Zahlen
\newcommand{\R}       {\mathbb{R}}          % reelle Zahlen
\newcommand{\Prob}    {\mathrm{P}}          % Wahrscheinlichkeit
\newcommand{\inter}   {\cap}                % Schnittmenge
\newcommand{\union}   {\cup}                % Vereinigung
\newcommand{\Oh}      {O}                   % O-Notation (Landau-Symbole)
\newcommand{\mycite}[1]{\textcolor{tumgreen}{[#1]}} 

\newenvironment{changemargin}[2]{% 
  \begin{list}{}{%
    \setlength{\topsep}{0pt}%
    \setlength{\leftmargin}{#1}%
    \setlength{\rightmargin}{#2}%
    \setlength{\listparindent}{\parindent}%
    \setlength{\itemindent}{\parindent}%
    \setlength{\parsep}{\parskip}%
  }%
  \item[]}{\end{list}}

\title{A flexible tool for Shape Analysis}
\subtitle{IDP Presentation}
\author{\href{emanuel.laude@in.tum.de}{Emanuel Laude}\\ \href{laehner@in.tum.de}{Zorah L\"ahner} } 
%\date{\today}
\date{January 16, 2015}
\institute{Technische Universit�t M�nchen}

% Inhaltsverzeichnis zu Begin von jedem Abschnitt einblenden?
%\AtBeginSection[]{
%  \begin{frame}
%    \frametitle{Outline}
%    \tableofcontents[currentsection]
%  \end{frame}
%}

\begin{document}

\begin{frame}
  \titlepage
\end{frame}

% Inhaltsverzeichnis
%\begin{frame}
%  \frametitle{Outline}
%  \tableofcontents
%\end{frame}

\section{Motivation}

\begin{frame}
  \frametitle{Motivation}
  \small
  \begin{block}{Observation 1}<1->
    \footnotesize
    Analyzing biological sequences for their statistical properties can lead to interesting conclusions:
    \begin{itemize}
      \footnotesize
      \item Coding vs. non-coding regions (CpG content, entropy, \ldots)
      \item Species identification
      \item Number and length of repeats
    \end{itemize}
  \end{block}
  \begin{block}{Observation 2}<2->
    \footnotesize
    Many algorithms and data structures depend on statistical properties of the underlying sequence:
    \begin{itemize}
      \footnotesize
      \item Algorithms for sequence assembly: repeats
      \item Suffix tree: repeats
      \item $q$-gram index: number and distribution of $q$-grams
      \item Compressed index structures: entropy
    \end{itemize}
  \end{block}
\end{frame}

\section{Algorithm}

\begin{frame}
  \frametitle{Logo}
  The logo of the Technische Universit�t M�nchen looks as follows:
  \only<2>{\includegraphics[width=\textwidth]{TU_Muenchen_Logo_Breit}}
\end{frame}

\section{Discussion}

\begin{frame}
  \frametitle{Discussion}
  \begin{itemize}
    \item<1-> Summary
      \begin{itemize}
        \item First
        \item Second
      \end{itemize}
    \item<2-> Outlook
      \begin{itemize}
        \item First
        \item Second
      \end{itemize}
  \end{itemize}
\end{frame}

\begin{frame}
  \Large
  \begin{itemize}
    \item<1> Thank you!
  \end{itemize}
  \begin{itemize}
    \item<2> Questions?
  \end{itemize}
\end{frame}

\end{document}
